% Four things `\halign` does that nothing in plain.tex or the corpus exercises,
% and that amsmath's `align` depends on all four of at once.
%
% 1. A *repeating* preamble. `tex.web` §783: an `&` arriving before any `#`,
%    with the template so far empty and no loop yet recorded, sets `cur_loop`
%    — the columns from there on repeat for as long as the rows need them
%    (§793 "lengthen the preamble periodically"). `\halign{&<template>\cr}` is
%    the whole idiom, and `\align@preamble` opens with exactly that `&`.
%
% 2. `\everycr`. §774 inserts it once the preamble's `\cr` is read, and §799
%    after every finished row. amsmath resets its per-row counters there, so
%    without it the field count accumulated across rows and the alignment came
%    out with as many spurious empty columns as it had rows.
%
% 3. `\setbox<n>=\lastbox` when there is no last box *voids* box n (§1079
%    `box_end` assigns the null box). Loops are built on this: amsmath measures
%    a row by peeling boxes off it until `\ifhbox` says there are none left, so
%    leaving the register alone instead means the loop never ends.
%
% 4. Token-list *parameters* are eqtb entries, so a group restores them. Only
%    the numbered `\toks` registers were getting that here, which made every
%    `{\everypar{…}}` permanent — and, since `\everycr` is one of them, made
%    one alignment's `\everycr` outlive it.
\catcode`\{=1 \catcode`\}=2 \catcode`\#=6 \catcode`\&=4
\font\testfont=cmr10 \testfont
\hsize=300pt \vsize=300pt \parindent=0pt \tabskip=3pt

% --- 1. the repeating preamble ---------------------------------------------
% Three columns from a one-column template, and a row that uses fewer.
\setbox0\vbox{\halign{&\hfil#\hfil\cr A&B&C\cr D&E\cr}}
\message{<loop wd=\the\wd0 ht=\the\ht0>}
% The same alignment with the columns written out must measure the same.
\setbox1\vbox{\halign{\hfil#\hfil&\hfil#\hfil&\hfil#\hfil\cr A&B&C\cr D&E\cr}}
\message{<fixed wd=\the\wd1 ht=\the\ht1>}
% A loop that starts partway through: column one is its own, the rest repeat.
\setbox2\vbox{\halign{#&&[#]\cr A&B&C\cr}}
\message{<partial wd=\the\wd2 ht=\the\ht2>}

% --- 2. \everycr ------------------------------------------------------------
% Once after the preamble, once after each row: three for two rows.
\count0=0
\setbox3\vbox{\everycr{\noalign{\global\advance\count0 by 1}}%
  \halign{#&#\cr A&B\cr C&D\cr}}
\message{<everycr=\the\count0>}
% It fires the same for a `\cr` that arrives from a macro, and for `\crcr`.
\count1=0
\def\myrow{A&B\cr}
\setbox3\vbox{\everycr{\noalign{\global\advance\count1 by 1}}%
  \halign{&#\crcr \myrow \myrow}}
\message{<everycr-macro=\the\count1>}

% --- 3. \lastbox with nothing to take --------------------------------------
\setbox4\hbox{A}
\setbox5\hbox{}
\setbox5\hbox{\unhbox5 \global\setbox4\lastbox}
\message{<lastbox-empty void=\ifvoid4 Y\else N\fi>}
\setbox4\hbox{A}
\setbox5\hbox{\kern3pt}
\setbox5\hbox{\unhbox5 \global\setbox4\lastbox}
\message{<lastbox-kern void=\ifvoid4 Y\else N\fi>}
% …and when there *is* one, it is taken off the list, so the box shrinks.
\setbox4\hbox{A}
\setbox5\hbox{\hbox{B}}
\setbox5\hbox{\unhbox5 \global\setbox4\lastbox}
\message{<lastbox-box void=\ifvoid4 Y\else N\fi wd=\the\wd5>}

% --- 4. token-list parameters are group-scoped ------------------------------
\everypar{A}
{\everypar{B}\message{<toks-inner=[\detokenize\expandafter{\the\everypar}]>}}
\message{<toks-after=[\detokenize\expandafter{\the\everypar}]>}
\begingroup\everypar{C}\endgroup
\message{<toks-after2=[\detokenize\expandafter{\the\everypar}]>}
% `\global` still leaks out on purpose.
{\global\everypar{D}}
\message{<toks-global=[\detokenize\expandafter{\the\everypar}]>}

\end
